% =============================================================================
% Implementação de Baixo Nível da Biblioteca CKKS para Aprendizado Federado
% Documento standalone – estilo dissertação (PT-BR)
% Compilar: pdflatex dissertacao_baixo_nivel.tex  (2 passagens)
% =============================================================================
\documentclass[11pt,a4paper]{report}

% --- Codificação e idioma ---
\usepackage[utf8]{inputenc}
\usepackage[T1]{fontenc}
\usepackage[brazil]{babel}

% --- Matemática ---
\usepackage{amsmath}
\usepackage{amssymb}
\usepackage{amsthm}

% --- Algoritmos ---
\usepackage{algorithm}
\usepackage{algpseudocode}
\floatname{algorithm}{Algoritmo}
\renewcommand{\listalgorithmname}{Lista de Algoritmos}

% --- Diagramas TikZ ---
% Nota: utiliza-se TikZ diretamente (bibliotecas arrows.meta, positioning, calc,
% decorations.pathreplacing). A família tkz-euclide/tkz-graph destina-se a
% construções de geometria euclidiana e teoria dos grafos, não sendo adequada
% para diagramas de borboleta, fluxo de sinal e pipeline presentes neste documento.
\usepackage{tikz}
\usetikzlibrary{arrows.meta, positioning, calc, decorations.pathreplacing, matrix, shapes.geometric}

% --- Tabelas ---
\usepackage{booktabs}

% --- Página ---
\usepackage[a4paper, top=3cm, bottom=2.5cm, left=3cm, right=2.5cm]{geometry}

% --- Links e referências ---
\usepackage[colorlinks=true, allcolors=blue,
            pdfauthor={Yuri Dimitre Dias de Faria},
            pdftitle={Implementação de Baixo Nível da Biblioteca CKKS}]{hyperref}
\usepackage[brazilian, nameinlink]{cleveref}

% --- Teoremas ---
\newtheorem{definition}{Definição}[chapter]
\newtheorem{proposition}{Proposição}[chapter]
\newtheorem{remark}{Observação}[chapter]

% --- Macros ---
\newcommand{\Rq}{\mathcal{R}_q}
\newcommand{\Zq}{\mathbb{Z}_q}
\newcommand{\npoly}{N_{\text{poly}}}

% =============================================================================
\begin{document}

\title{\textbf{Implementação de Baixo Nível da Biblioteca CKKS\\
para Aprendizado Federado Preservador de Privacidade}\\[1em]
\large Documento técnico complementar à dissertação de mestrado}
\author{Yuri Dimitre Dias de Faria\\
\small Instituto de Computação -- UNICAMP\\
\small Orientador: Prof. Dr. Allan Mariano de Souza}
\date{Junho de 2026}
\maketitle

\tableofcontents
\listoffigures
\listofalgorithms

% =============================================================================
\chapter{Introdução}
\label{cap:intro}
% =============================================================================

Este documento descreve a arquitetura e os mecanismos de baixo nível do \emph{backend}
C++ da biblioteca CKKS customizada, desenvolvida para suportar operações de cifragem
homomórfica em cenários de Aprendizado Federado (FL). O material aqui apresentado
é complementar ao Capítulo~4 da dissertação de mestrado associada, que cobre os fundamentos
do esquema CKKS, o modelo de erro da agregação homomórfica e a parametrização do
problema RLWE~\cite{Cheon2017,sealmanual}. Pressupõe-se familiaridade com esses
conceitos; o foco deste documento é a tradução das estruturas matemáticas em código
executável de alto desempenho.

O núcleo da biblioteca, denominado \texttt{ckks\_core}, é implementado em C++17 e
compilado como biblioteca estática com as diretivas de otimização
\texttt{-O3 -march=native -ffast-math -fopenmp-simd}, que ativam vetorização SIMD
sem recorrer ao \emph{runtime} de \emph{threads} OpenMP, cuja integração com os
laços da NTT está prevista para uma versão futura. A \cref{tab:modulos} lista os
seis módulos hierárquicos e os arquivos correspondentes.

\begin{table}[ht]
\centering
\caption{Módulos do núcleo \texttt{ckks\_core} e arquivos fonte.}
\label{tab:modulos}
\begin{tabular}{llp{7cm}}
\toprule
\textbf{Módulo} & \textbf{Arquivos} & \textbf{Responsabilidade} \\
\midrule
1 – Aritmética modular & \texttt{math\_utils.h/.cpp} & \texttt{mod\_add}, \texttt{mod\_sub}, Montgomery, primos NTT \\
2 – Polinômio & \texttt{poly.h/.cpp} & Tipo \texttt{Poly} e adição SIMD \\
3 – NTT & \texttt{ntt.h/.cpp} & Transformada negacíclica $O(N\log N)$ \\
4 – Amostragem & \texttt{sampler.h/.cpp} & Uniforme, gaussiana discreta, ternária \\
5 – Chave/Cifragem & \texttt{keygen.h}, \texttt{encryptor.h}, \texttt{decryptor.h}, \texttt{ckks\_core.cpp} & Geração de chaves, \texttt{encrypt}, \texttt{decrypt} \\
6 – Avaliação/Codificação & \texttt{evaluator.h/.cpp}, \texttt{encoder.h/.cpp}, \texttt{serialization.cpp} & Adição homomórfica, codificação canônica, serialização \\
\bottomrule
\end{tabular}
\end{table}

Em contraste com implementações de propósito geral como Microsoft SEAL~\cite{sealmanual},
o \texttt{ckks\_core} adota as simplificações a seguir, motivadas pelo foco exclusivo
em agregação federada por adição homomórfica:

\begin{itemize}
  \item \textbf{Módulo único:} utiliza um único primo NTT $q \equiv 1 \pmod{2N}$, com
    $q < 2^{62}$, em lugar da cadeia $Q = \prod_j q_j$ do modelo RNS completo discutido
    no Capítulo~4 da dissertação. Não há \emph{rescaling} nem consumo progressivo de níveis.
  \item \textbf{Regime somente-aditivo:} o servidor de FL agrega gradientes apenas por
    $\texttt{Ciphertext} + \texttt{Ciphertext}$; não são implementadas relinearização,
    chaves de Galois nem \emph{bootstrapping}.
  \item \textbf{Orçamento de ruído estável:} a adição homomórfica não consome níveis da
    cadeia de módulos; o orçamento de ruído permanece constante ao longo das rodadas de
    agregação, ao contrário dos circuitos multiplicativos profundos.
\end{itemize}

Essas escolhas permitem concentrar todo o custo computacional nas multiplicações
polinomiais da NTT, que dominam a cifragem e a decifragem.

% =============================================================================
\chapter{Fundamentos Algébricos do Anel de Polinômios}
\label{cap:ring}
% =============================================================================

A biblioteca opera sobre o anel quociente
\begin{equation}
  \Rq \;=\; \mathbb{Z}_q[X]\,/\,(X^N + 1),
  \label{eq:ring}
\end{equation}
onde $q$ é um número primo, $N$ é uma potência de dois e os coeficientes pertencem a
$\Zq = \{0, 1, \ldots, q-1\}$. A relação $X^N \equiv -1 \pmod{X^N+1}$ confere ao anel
uma estrutura negacíclica: no produto de dois polinômios $a, b \in \Rq$, o coeficiente
de índice $i+j \geq N$ contribui com sinal negativo para o coeficiente de índice $i+j-N$.

Na notação da dissertação, o grau polinomial aqui denominado $N$ corresponde ao parâmetro
$\npoly$. A distinção fundamental deste \emph{backend} é que o módulo $q$ é um único
primo, ao contrário da cadeia $Q = \prod_j q_j$ do modelo RNS; não há portanto
\emph{rescaling} nem índice de nível.

\begin{definition}[Primo NTT-amigável]
\label{def:ntt_prime}
Um primo $q$ é dito NTT-amigável de parâmetro $N$ se $q \equiv 1 \pmod{2N}$. Esta
congruência garante a existência de uma raiz primitiva $2N$-ésima da unidade
$\psi \in \mathbb{Z}_q^*$ tal que $\psi^N \equiv -1 \pmod{q}$, necessária para a
NTT negacíclica descrita no \cref{cap:ntt}.
\end{definition}

Em memória, um polinômio é representado pela estrutura \texttt{Poly} com três campos:
o módulo \texttt{uint64\_t q}, o grau \texttt{uint32\_t N} e um vetor de $N$ coeficientes
\texttt{std::vector<uint64\_t> coeffs} em $[0, q)$. A adição de dois polinômios é
implementada coeficiente a coeficiente com a diretiva \texttt{\#pragma omp simd}, que
instrui o compilador a emitir instruções AVX2 (\texttt{vpsubq}, \texttt{vpcmpgtq},
\texttt{vpand}) quando compilado com \texttt{-march=native}. A redução modular em cada
coeficiente é realizada pela função \texttt{mod\_add}, descrita no \cref{cap:modular}.

% =============================================================================
\chapter{Aritmética Modular Sem Desvios}
\label{cap:modular}
% =============================================================================

Adição e subtração modulares são os blocos primitivos de todos os laços internos da
biblioteca. Para que o compilador os vetorize de forma eficiente, é necessário que não
contenham desvios condicionais (\emph{branches}). A implementação emprega o truque da
máscara binária: uma comparação de inteiros de 64 bits produz um valor em $\{0, 1\}$
cuja negação bitwise em aritmética sem sinal produz a máscara $\{0,\; 2^{64}-1\}$.

\section{Adição modular}
\label{sec:mod_add}

Sejam $a, b \in [0, q)$ com $q < 2^{63}$. A soma $a + b \in [0, 2q)$, de modo que é
necessário subtrair $q$ se e somente se $a + b \geq q$. O \cref{alg:mod_add} implementa
esta correção sem desvio.

\begin{algorithm}[ht]
\caption{\texttt{mod\_add}: adição modular sem desvio condicional}
\label{alg:mod_add}
\begin{algorithmic}[1]
\Require $a, b \in [0, q)$;\; $q < 2^{63}$
\Ensure $r = (a + b) \bmod q$
\State $c \gets a + b$
  \Comment{Sem \emph{overflow}: $c \in [0, 2q) \subset [0, 2^{64})$}
\State $\text{máscara} \gets -(c \geq q)$
  \Comment{$0$ se $c < q$;\; $2^{64}-1$ se $c \geq q$}
\State \Return $c - (q \;\mathbin{\&}\; \text{máscara})$
\end{algorithmic}
\end{algorithm}

A expressão $-(c \geq q)$, em aritmética inteira de 64 bits sem sinal, produz $0$ quando
$c < q$ e $2^{64}-1$ quando $c \geq q$. O operador \texttt{\&} com $q$ seleciona
$0$ ou $q$ sem desvio; o resultado pertence a $[0, q)$.

\section{Subtração modular}
\label{sec:mod_sub}

Sejam $a, b \in [0, q)$. Se $a \geq b$, o resultado é $a - b \in [0, q)$; se $a < b$,
o valor correto é $a - b + q$. O \cref{alg:mod_sub} aplica a correção pela mesma
técnica de máscara.

\begin{algorithm}[ht]
\caption{\texttt{mod\_sub}: subtração modular sem desvio condicional}
\label{alg:mod_sub}
\begin{algorithmic}[1]
\Require $a, b \in [0, q)$;\; $q < 2^{63}$
\Ensure $r = (a - b) \bmod q$
\State $c \gets a - b$
  \Comment{Enrola para $[2^{64}-q,\, 2^{64})$ se $a < b$}
\State $\text{máscara} \gets -(a < b)$
\State \Return $c + (q \;\mathbin{\&}\; \text{máscara})$
\end{algorithmic}
\end{algorithm}

As duas funções são declaradas \texttt{inline noexcept} no cabeçalho
\texttt{include/math\_utils.h} e produzem, sob \texttt{-O3 -march=native}, código
que o compilador reduz a duas a três instruções AVX2 vetoriais por laço de $N$
coeficientes, eliminando qualquer desvio condicional no caminho crítico.

% =============================================================================
\chapter{Redução de Montgomery}
\label{cap:montgomery}
% =============================================================================

A multiplicação modular ingênua requer uma divisão por $q$ a cada produto, operação
custosa em hardware de propósito geral. A redução de Montgomery~\cite{Montgomery1985}
substitui essa divisão por operações de deslocamento e soma, representando os
elementos de $\mathbb{Z}_q$ em uma forma que incorpora um fator de potência de dois $R$.

\section{Forma de Montgomery e constantes pré-computadas}
\label{sec:mont_constants}

Adota-se $R = 2^{64}$ e exige-se $q < 2^{62}$ (primo ímpar). A forma de Montgomery de
$a \in \mathbb{Z}_q$ é
\begin{equation}
  \hat{a} \;=\; a \cdot R \bmod q.
  \label{eq:mont_form}
\end{equation}
A estrutura \texttt{MontgomeryCtx} armazena três constantes pré-computadas em
tempo de construção:
\begin{itemize}
  \item $q_{\text{inv}} = -q^{-1} \bmod 2^{64}$: usado internamente pelo algoritmo REDC;
  \item $r_2 = R^2 \bmod q$: permite converter $a$ para a forma de Montgomery via
    $\hat{a} = \text{REDC}(a \cdot r_2)$;
  \item $q$: o módulo primo.
\end{itemize}

O cálculo de $q_{\text{inv}}$ emprega iteração de Newton sobre $\mathbb{Z}/2^{64}\mathbb{Z}$,
que dobra a precisão binária a cada passo, conforme o \cref{alg:qinv}.

\begin{algorithm}[ht]
\caption{Cálculo de $q_{\text{inv}} = -q^{-1} \bmod 2^{64}$ por Newton--Raphson}
\label{alg:qinv}
\begin{algorithmic}[1]
\Require $q$ ímpar, $q < 2^{62}$
\Ensure $q_{\text{inv}}$ tal que $q \cdot q_{\text{inv}} \equiv -1 \pmod{2^{64}}$
\State $x \gets q$
  \Comment{Semente: $q \cdot q \equiv 1 \pmod{8}$ para todo $q$ ímpar}
\For{$i \gets 1$ \textbf{até} $5$}
  \State $x \gets x \cdot (2 - q \cdot x)$
    \Comment{Dobra a precisão: de $\sim 3$ para $\sim 96$ bits em 5 iterações}
\EndFor
\State \Return $-x \bmod 2^{64}$
\end{algorithmic}
\end{algorithm}

O valor $r_2$ é obtido via aritmética de 128 bits:
\begin{equation}
  r_2 = (R \bmod q)^2 \bmod q.
  \label{eq:r2}
\end{equation}

\section{Algoritmo REDC}
\label{sec:redc}

O algoritmo REDC~\cite{Montgomery1985} computa $T \cdot R^{-1} \bmod q$ para
$T \in [0, q \cdot R)$ sem divisão explícita por $q$.

\begin{algorithm}[ht]
\caption{REDC: redução de Montgomery}
\label{alg:redc}
\begin{algorithmic}[1]
\Require $T \in [0, q \cdot R)$;\; $q_{\text{inv}} = -q^{-1} \bmod R$;\; $R = 2^{64}$
\Ensure $T \cdot R^{-1} \bmod q$
\State $m \gets (T \bmod R) \cdot q_{\text{inv}} \bmod R$
  \Comment{Escolhe $m$ tal que $T + mq \equiv 0 \pmod{R}$}
\State $t \gets (T + m \cdot q) \div R$
  \Comment{Divisão exata: descarta os 64 bits inferiores por deslocamento}
\State \Return $t - (q \;\mathbin{\&}\; -(t \geq q))$
  \Comment{Normaliza $t \in [0, 2q)$ para $[0, q)$ sem desvio}
\end{algorithmic}
\end{algorithm}

A propriedade $T + mq \equiv 0 \pmod{R}$ garante que a divisão na linha~2 seja exata,
implementável como um deslocamento de 64 bits. O resultado preliminar pertence a $[0, 2q)$;
a normalização da linha~3 usa o mesmo truque de máscara da \cref{sec:mod_add}.

\section{Operações públicas do contexto de Montgomery}
\label{sec:mont_ops}

As três operações do \texttt{MontgomeryCtx} são definidas em termos de REDC:
\begin{align}
  \texttt{to\_mont}(a) &= \text{REDC}(a \cdot r_2) = a \cdot R \bmod q = \hat{a},
    \label{eq:to_mont}\\
  \texttt{from\_mont}(\hat{a}) &= \text{REDC}(\hat{a}) = a,
    \label{eq:from_mont}\\
  \texttt{mont\_mul}(\hat{a}, \hat{b}) &= \text{REDC}(\hat{a} \cdot \hat{b})
    = a \cdot b \cdot R \bmod q = \widehat{a \cdot b}.
    \label{eq:mont_mul}
\end{align}

A \cref{fig:montgomery_flow} ilustra o padrão de uso: os operandos são convertidos
para a forma de Montgomery uma única vez; todas as multiplicações intermediárias operam
no domínio de Montgomery sem custo adicional de conversão; a conversão inversa ocorre
apenas ao final.

\begin{figure}[ht]
\centering
\begin{tikzpicture}[
  node distance = 0.6cm and 1.5cm,
  box/.style  = {draw, rounded corners=3pt, minimum width=2.6cm, minimum height=0.65cm,
                 align=center, font=\small},
  io/.style   = {draw, ellipse, minimum width=2.0cm, minimum height=0.65cm,
                 align=center, font=\small},
  arr/.style  = {-{Stealth[length=5pt]}, thick}
]
  \node[io]  (a)    {$a \in \mathbb{Z}_q$};
  \node[io,  below=0.7cm of a] (b) {$b \in \mathbb{Z}_q$};

  \node[box, right=1.4cm of a]  (am)   {$\hat{a} = aR \bmod q$};
  \node[box, right=1.4cm of b]  (bm)   {$\hat{b} = bR \bmod q$};

  \node[box, right=1.4cm of am, yshift=-0.68cm] (mul) {\texttt{mont\_mul}\\$\hat{a}\cdot\hat{b}$};

  \node[box, right=1.4cm of mul] (res)  {$\widehat{ab}$};
  \node[io,  right=1.3cm of res] (out)  {$ab \in \mathbb{Z}_q$};

  \draw[arr] (a)  -- node[above, font=\footnotesize]{\texttt{to\_mont}} (am);
  \draw[arr] (b)  -- node[below, font=\footnotesize]{\texttt{to\_mont}} (bm);
  \draw[arr] (am) |- (mul);
  \draw[arr] (bm) |- (mul);
  \draw[arr] (mul) -- node[above, font=\footnotesize]{REDC} (res);
  \draw[arr] (res) -- node[above, font=\footnotesize]{\texttt{from\_mont}} (out);
\end{tikzpicture}
\caption{Fluxo de uso típico do domínio de Montgomery. As conversões ocorrem uma única vez;
o laço de multiplicações intermediárias usa exclusivamente \texttt{mont\_mul}, sem divisão
por $q$.}
\label{fig:montgomery_flow}
\end{figure}

% =============================================================================
\chapter{Geração de Primos NTT-Amigáveis}
\label{cap:primes}
% =============================================================================

O módulo primo $q$ é gerado pela função \texttt{generate\_proth\_prime(bit\_size, N)},
que encontra o maior primo da forma $q = 2Nk + 1$ com exatamente \texttt{bit\_size} bits,
satisfazendo assim a congruência $q \equiv 1 \pmod{2N}$ da \cref{def:ntt_prime} e
o limite $q < 2^{62}$ exigido pelo contexto de Montgomery.

\begin{algorithm}[ht]
\caption{\texttt{generate\_proth\_prime}: primo NTT-amigável de tamanho fixo}
\label{alg:prime_gen}
\begin{algorithmic}[1]
\Require $\texttt{bit\_size} \in [2, 62]$;\; $N > 0$
\Ensure Primo $q \equiv 1 \pmod{2N}$,\; $2^{\texttt{bit\_size}-1} \leq q < 2^{\texttt{bit\_size}}$
\State $\text{passo} \gets 2N$
\State $q \gets \lfloor (2^{\texttt{bit\_size}} - 2) / \text{passo} \rfloor
       \cdot \text{passo} + 1$
  \Comment{Maior candidato $< 2^{\texttt{bit\_size}}$ com $q \equiv 1 \pmod{2N}$}
\While{$q \geq 2^{\texttt{bit\_size}-1}$}
  \If{\Call{MillerRabin}{$q$}}
    \State \Return $q$
  \EndIf
  \State $q \gets q - \text{passo}$
\EndWhile
\State \textbf{erro:} nenhum primo encontrado para os parâmetros fornecidos
\end{algorithmic}
\end{algorithm}

O teste de primalidade \textsc{MillerRabin} é a variante determinística com as 12
testemunhas fixas $\{2, 3, 5, 7, 11, 13, 17, 19, 23, 29, 31, 37\}$, cuja exatidão
é garantida para todo $n < 3{,}3 \times 10^{24}$, cobrindo o intervalo $[0, 2^{62})$
sem falsos positivos~\cite{Pomerance1980}. Cada candidato satisfaz
$q \equiv 1 \pmod{2N}$ por construção; a busca em ordem descendente seleciona o primo
de maior valor com a largura de bits solicitada, maximizando a margem disponível para o
orçamento de ruído $B = \log_2 q$.

% =============================================================================
\chapter{Transformada Numérica Teórica Negacíclica}
\label{cap:ntt}
% =============================================================================

A multiplicação de polinômios em $\Rq$ constitui a operação dominante tanto na cifragem
quanto na decifragem. A multiplicação escolar tem complexidade $O(N^2)$; a Transformada
Numérica Teórica (NTT) reduz esse custo para $O(N \log N)$~\cite{LongaNaehrig2016}.
A implementação utiliza a variante negacíclica, que avalia os polinômios nas raízes
primitivas $2N$-ésimas da unidade, tornando a redução módulo $X^N+1$ implícita ao
próprio algoritmo.

\section{Derivação da raiz primitiva $\psi$}
\label{sec:psi}

Dado $q \equiv 1 \pmod{2N}$, o grupo multiplicativo $\mathbb{Z}_q^*$ admite uma raiz
primitiva $g$ (menor inteiro cujas potências percorrem todo o grupo). O elemento
\begin{equation}
  \psi = g^{(q-1)/(2N)} \bmod q
  \label{eq:psi}
\end{equation}
tem ordem exatamente $2N$, satisfazendo $\psi^N \equiv -1 \pmod{q}$. Esta verificação
é realizada no construtor do contexto \texttt{NttCtx}; caso contrário, uma exceção é
lançada, prevenindo uso silencioso de parâmetros inválidos.

\section{Pré-computação das tabelas de twiddle}
\label{sec:twiddle}

O contexto \texttt{NttCtx} armazena, além de $\psi$ e $\psi^{-1}$, dois vetores de
$N$ elementos chamados \emph{tabelas de twiddle}:
\begin{align}
  \texttt{psi\_pow}[k] &= \texttt{to\_mont}(\psi^{\,\text{bit\_rev}(k)}),
    \label{eq:psi_pow}\\
  \texttt{psi\_inv\_pow}[k] &= \texttt{to\_mont}(\psi^{-\text{bit\_rev}(k)}),
    \label{eq:psi_inv_pow}
\end{align}
onde $\text{bit\_rev}(k)$ é o inteiro obtido invertendo os $\log_2 N$ bits de $k$.

As potências são geradas por multiplicação sequencial de Montgomery, conforme a
\cref{eq:twiddle_seq}, evitando a exponenciação por quadração para cada entrada:
\begin{equation}
  \texttt{raw}[k] = \psi^k \bmod q
    = \texttt{mont\_mul}(\texttt{raw}[k-1],\; \hat{\psi}).
  \label{eq:twiddle_seq}
\end{equation}
Em seguida, os valores são permutados pela tabela \texttt{bit\_rev}, de modo que o
laço interno da borboleta acessa \texttt{psi\_pow} em sequência linear, favorecendo
a linha de \emph{cache}. Adicionalmente, o contexto armazena
\begin{equation}
  N_{\text{inv}}^{\text{mont}} = \texttt{to\_mont}(N^{-1} \bmod q),
  \label{eq:ninv}
\end{equation}
utilizado para normalizar a NTT inversa.

\section{NTT direta: Cooley--Tukey DIT}
\label{sec:ntt_forward}

A NTT direta (\emph{decimation-in-time}, DIT) do \cref{alg:ntt_forward} processa
$\log_2 N$ estágios em ordem crescente de $m$. Em cada estágio, borboletas radix-2
operam sobre pares de coeficientes separados por \texttt{gap} posições.

\begin{algorithm}[ht]
\caption{\texttt{ntt\_forward}: NTT negacíclica direta (Cooley--Tukey DIT)}
\label{alg:ntt_forward}
\begin{algorithmic}[1]
\Require Vetor $a[0\ldots N-1]$ em $\mathbb{Z}_q$;\; contexto \texttt{NttCtx}
\Ensure $a$ substituído pelos valores de avaliação nas raízes $\psi^1, \psi^3, \ldots, \psi^{2N-1}$
\For{$m \gets 1$ \textbf{até} $N/2$ por duplicação}
  \State $\text{gap} \gets N / (2m)$
  \For{$i \gets 0$ \textbf{até} $m-1$}
    \State $w_i \gets \texttt{psi\_pow}[m + i]$
      \Comment{Fator de giro em forma de Montgomery}
    \For{$j \gets 0$ \textbf{até} $\text{gap}-1$}
      \State $u \gets a[2i \cdot \text{gap} + j]$
      \State $v \gets \texttt{mont\_mul}(a[2i \cdot \text{gap} + \text{gap} + j],\; w_i)$
        \Comment{Twiddle \emph{antes} da soma/diferença (DIT)}
      \State $a[2i \cdot \text{gap} + j]            \gets \texttt{mod\_add}(u, v, q)$
      \State $a[2i \cdot \text{gap} + \text{gap} + j] \gets \texttt{mod\_sub}(u, v, q)$
    \EndFor
  \EndFor
\EndFor
\end{algorithmic}
\end{algorithm}

Entrada em ordem natural; saída em ordem bit-reversa. Os ponteiros \texttt{\_\_restrict\_\_}
no laço interno informam ao compilador a ausência de aliasing entre os vetores, favorecendo
a vetorização automática.

\section{NTT inversa: Gentleman--Sande DIF}
\label{sec:ntt_inverse}

A NTT inversa (\emph{decimation-in-frequency}, DIF) é a transposta algébrica da direta.
Ela processa os estágios em ordem decrescente de $m$, realizando primeiro a soma/diferença
e depois a multiplicação pelo fator de giro inverso.

\begin{algorithm}[ht]
\caption{\texttt{ntt\_inverse}: NTT negacíclica inversa (Gentleman--Sande DIF)}
\label{alg:ntt_inverse}
\begin{algorithmic}[1]
\Require Vetor $\hat{a}[0\ldots N-1]$ em ordem bit-reversa;\; contexto \texttt{NttCtx}
\Ensure $a$ em ordem natural (coeficientes do polinômio original)
\For{$m \gets N/2$ \textbf{descendo até} $1$ por divisão por 2}
  \State $\text{gap} \gets N / (2m)$
  \For{$i \gets 0$ \textbf{até} $m-1$}
    \State $w_i \gets \texttt{psi\_inv\_pow}[m + i]$
    \For{$j \gets 0$ \textbf{até} $\text{gap}-1$}
      \State $u \gets a[2i \cdot \text{gap} + j]$
      \State $v \gets a[2i \cdot \text{gap} + \text{gap} + j]$
      \State $a[2i \cdot \text{gap} + j]            \gets \texttt{mod\_add}(u, v, q)$
        \Comment{Soma \emph{antes} do twiddle (DIF)}
      \State $a[2i \cdot \text{gap} + \text{gap} + j] \gets
        \texttt{mont\_mul}(\texttt{mod\_sub}(u, v, q),\; w_i)$
    \EndFor
  \EndFor
\EndFor
\For{$i \gets 0$ \textbf{até} $N-1$}
  \State $a[i] \gets \texttt{mont\_mul}(a[i],\; N_{\text{inv}}^{\text{mont}})$
    \Comment{Normalização: remove fator $N$ acumulado}
\EndFor
\end{algorithmic}
\end{algorithm}

Após $\log_2 N$ estágios DIF, cada coeficiente carrega um fator implícito $N$, removido
pela passagem de normalização final. Os algoritmos \cref{alg:ntt_forward} e
\cref{alg:ntt_inverse} formam par exato: $\texttt{ntt\_inverse}(\texttt{ntt\_forward}(a)) = a$
para qualquer $a \in \Rq$, sem permutação adicional.

\section{Diagramas da borboleta e do grafo de fluxo de sinal}
\label{sec:ntt_diagrams}

A \cref{fig:butterfly} apresenta as duas variantes da borboleta radix-2. A
\cref{fig:ntt8} mostra o grafo de fluxo de sinal completo da NTT direta para $N=8$,
com três estágios.

\begin{figure}[ht]
\centering
\begin{tikzpicture}[
  arr/.style = {-{Stealth[length=4pt]}, thick},
  neg/.style = {-{Stealth[length=4pt]}, thick, dashed},
  lbl/.style = {font=\small}
]
%% --- Borboleta DIT (esquerda) ---
\begin{scope}
  % nós de entrada
  \node[lbl] (ul) at (0,  1.0) {$u$};
  \node[lbl] (vl) at (0,  0.0) {$v$};
  % nós de saída
  \node[lbl] (ur) at (3.8, 1.0) {$u + w{\cdot}v$};
  \node[lbl] (vr) at (3.8, 0.0) {$u - w{\cdot}v$};
  % setas
  \draw[arr] (0.2, 1.0) -- (2.5, 1.0);
  \draw[arr] (0.2, 0.0) -- (2.5, 0.0);
  \draw[arr] (0.2, 0.0) -- node[right=1pt, font=\footnotesize]{$\times w$} (2.4, 0.92);
  \draw[neg] (0.2, 1.0) -- (2.4, 0.08);
  \node[lbl] at (1.9, -0.55) {(a) DIT (\emph{Cooley--Tukey})};
\end{scope}

%% --- Borboleta DIF (direita) ---
\begin{scope}[xshift=6.2cm]
  \node[lbl] (ul2) at (0,  1.0) {$u$};
  \node[lbl] (vl2) at (0,  0.0) {$v$};
  \node[lbl] (ur2) at (3.6, 1.0) {$u + v$};
  \node[lbl] (vr2) at (3.6, 0.0) {$w{\cdot}(u - v)$};
  \draw[arr] (0.2, 1.0) -- (2.7, 1.0);
  \draw[arr] (0.2, 0.0) -- (2.7, 0.0);
  \draw[arr] (0.2, 1.0) -- (2.5, 0.08);
  \draw[neg] (0.2, 0.0) -- node[right=1pt, font=\footnotesize]{$\times w$} (2.5, 0.92);
  \node[lbl] at (1.8, -0.55) {(b) DIF (\emph{Gentleman--Sande})};
\end{scope}
\end{tikzpicture}
\caption{Borboletas radix-2. Setas sólidas indicam contribuição positiva; setas tracejadas
indicam contribuição negativa (subtração). (a)~DIT: o fator de giro $w$ é aplicado antes
da soma/diferença. (b)~DIF: o fator de giro é aplicado após a subtração.}
\label{fig:butterfly}
\end{figure}

\begin{figure}[ht]
\centering
\begin{tikzpicture}[
  xscale=2.0, yscale=0.72,
  dot/.style  = {circle, fill=black, inner sep=1.8pt},
  arr/.style  = {-{Stealth[length=3.5pt]}, thin},
  neg/.style  = {-{Stealth[length=3.5pt]}, thin, densely dashed},
  lbl/.style  = {font=\footnotesize}
]
% Estágios: 0=entrada, 1,2,3=saída
\foreach \stage in {0,1,2,3}{
  \foreach \row in {0,...,7}{
    \node[dot] (s\stage\row) at (\stage, -\row) {};
  }
}

% Rótulos de entrada
\foreach \i in {0,...,7}{
  \node[lbl, left=4pt of s0\i] {$x_{\i}$};
}

% Rótulos de saída (ordem bit-reversa: 0,4,2,6,1,5,3,7)
\foreach \i/\br in {0/0,1/4,2/2,3/6,4/1,5/5,6/3,7/7}{
  \node[lbl, right=4pt of s3\i] {$\hat{X}_{\br}$};
}

% Estágio 1: m=1, gap=4 — pares (0,4),(1,5),(2,6),(3,7)
\foreach \top/\bot in {0/4, 1/5, 2/6, 3/7}{
  \draw[arr] (s0\top) -- (s1\top);
  \draw[arr] (s0\bot) -- (s1\top);
  \draw[arr] (s0\top) -- (s1\bot);
  \draw[neg] (s0\bot) -- (s1\bot);
}

% Estágio 2: m=2, gap=2 — pares (0,2),(1,3),(4,6),(5,7)
\foreach \top/\bot in {0/2, 1/3, 4/6, 5/7}{
  \draw[arr] (s1\top) -- (s2\top);
  \draw[arr] (s1\bot) -- (s2\top);
  \draw[arr] (s1\top) -- (s2\bot);
  \draw[neg] (s1\bot) -- (s2\bot);
}

% Estágio 3: m=4, gap=1 — pares (0,1),(2,3),(4,5),(6,7)
\foreach \top/\bot in {0/1, 2/3, 4/5, 6/7}{
  \draw[arr] (s2\top) -- (s3\top);
  \draw[arr] (s2\bot) -- (s3\top);
  \draw[arr] (s2\top) -- (s3\bot);
  \draw[neg] (s2\bot) -- (s3\bot);
}

% Rótulos dos estágios
\node[lbl] at (0,  1.2) {Entrada};
\node[lbl] at (1,  1.2) {Estágio 1};
\node[lbl] at (2,  1.2) {Estágio 2};
\node[lbl] at (3,  1.2) {Estágio 3};
\end{tikzpicture}
\caption{Grafo de fluxo de sinal da NTT direta negacíclica para $N=8$ (3 estágios,
$\log_2 8 = 3$). Setas sólidas representam adição e setas tracejadas representam subtração
(fatores de giro omitidos por clareza). Entrada em ordem natural
$x_0, \ldots, x_7$; saída em ordem bit-reversa
$\hat{X}_0, \hat{X}_4, \hat{X}_2, \hat{X}_6, \hat{X}_1, \hat{X}_5, \hat{X}_3, \hat{X}_7$.}
\label{fig:ntt8}
\end{figure}

\section{Multiplicação polinomial via NTT}
\label{sec:poly_ntt_mul}

A multiplicação de $a, b \in \Rq$ é realizada em três etapas pelo \cref{alg:poly_ntt_mul},
com complexidade total $O(N \log N)$.

\begin{algorithm}[ht]
\caption{\texttt{poly\_ntt\_mul}: multiplicação polinomial negacíclica via NTT}
\label{alg:poly_ntt_mul}
\begin{algorithmic}[1]
\Require Polinômios $a, b \in \Rq$;\; contexto \texttt{NttCtx}
\Ensure $c = a \cdot b \bmod (X^N+1,\; q)$
\State $\hat{a} \gets \texttt{ntt\_forward}(a)$
\State $\hat{b} \gets \texttt{ntt\_forward}(b)$
\For{$i \gets 0$ \textbf{até} $N-1$}
  \State $\hat{c}[i] \gets (\hat{a}[i] \cdot \hat{b}[i]) \bmod q$
    \Comment{Produto de 128 bits; saídas da NTT são resíduos normais, não em forma de Montgomery}
\EndFor
\State \Return $\texttt{ntt\_inverse}(\hat{c})$
\end{algorithmic}
\end{algorithm}

O produto ponto a ponto utiliza aritmética de 128 bits em vez de \texttt{mont\_mul}
porque as avaliações produzidas por \texttt{ntt\_forward} são resíduos em forma normal
(os fatores de Montgomery dos twiddles se cancelam durante as borboletas), e a redução
direta por $q$ é computacionalmente equivalente, porém mais simples.

% =============================================================================
\chapter{Amostragem}
\label{cap:sampling}
% =============================================================================

O módulo de amostragem fornece as distribuições necessárias para a geração de chaves e
cifragem. A classe \texttt{Sampler} encapsula um gerador Mersenne Twister de 64 bits
(\texttt{mt19937\_64}~\cite{Matsumoto1998}), é não copiável para evitar correlação de
estados entre instâncias, e aceita semente determinística para reprodutibilidade em testes.

\section{Distribuição uniforme}
\label{sec:uniform}

A função \texttt{uniform(N, q)} gera um polinômio com cada coeficiente amostrado
independentemente e uniformemente de $\{0, 1, \ldots, q-1\}$, delegando a uniformidade
exata à distribuição \texttt{std::uniform\_int\_distribution<uint64\_t>(0, q-1)} da
biblioteca padrão C++.

\section{Gaussiana discreta}
\label{sec:gaussian}

A função \texttt{gaussian(N, q, sigma, tail)} amostra cada coeficiente segundo a
gaussiana discreta $D_{\mathbb{Z}, \sigma}$ com rejeição de amostras além de
$\text{tail} \cdot \sigma$ (padrão $\text{tail} = 6$, probabilidade de rejeição
$< 2 \times 10^{-9}$ por coeficiente), conforme o \cref{alg:gaussian}.

\begin{algorithm}[ht]
\caption{\texttt{gaussian}: amostragem da gaussiana discreta com rejeição}
\label{alg:gaussian}
\begin{algorithmic}[1]
\Require $N$, $q$, $\sigma > 0$, $\text{tail} \geq 1$
\Ensure Polinômio $e$ com coeficientes em $[0,q)$ distribuídos segundo $D_{\mathbb{Z},\sigma}$ truncada
\State $L \gets \lfloor \text{tail} \cdot \sigma \rceil$
\For{$i \gets 0$ \textbf{até} $N-1$}
  \Repeat
    \State $x \gets \lfloor \mathcal{N}(0, \sigma^2) \rceil$
      \Comment{Amostra normal e arredonda para inteiro}
  \Until{$|x| \leq L$}
  \State $r \gets x \bmod q$;\quad \textbf{se} $r < 0$ \textbf{então} $r \gets r + q$
    \Comment{Levanta $-1 \mapsto q-1$}
  \State $e[i] \gets r$
\EndFor
\State \Return $e$
\end{algorithmic}
\end{algorithm}

O parâmetro padrão $\sigma = 3{,}2$ segue a recomendação da literatura para o nível
de segurança de 128 bits no modelo RLWE~\cite{Cheon2017}. O ruído de cifragem gerado
por esta distribuição corresponde ao termo $e$ da análise de erro do Capítulo~4 da
dissertação.

\section{Distribuições ternárias}
\label{sec:ternary}

Dois tipos de distribuição ternária são suportados, ambos necessários em diferentes
pontos do protocolo.

\paragraph{Ternária uniforme.} \texttt{ternary\_uniform(N, q)} amostra cada coeficiente
independentemente de $\{-1, 0, 1\}$ com probabilidade $1/3$ cada, mapeando
$-1 \mapsto q-1$.

\paragraph{Ternária de peso de Hamming fixo.} \texttt{ternary\_hw(N, q, hw)} gera
exatamente $\text{hw}/2$ coeficientes iguais a $+1$ e $\text{hw}/2$ iguais a $q-1$,
com as demais posições nulas, usando o algoritmo de Fisher--Yates
parcial~\cite{Knuth1997}, conforme o \cref{alg:ternary_hw}.

\begin{algorithm}[ht]
\caption{\texttt{ternary\_hw}: ternário de peso de Hamming fixo}
\label{alg:ternary_hw}
\begin{algorithmic}[1]
\Require $N$, $q$, $\text{hw}$ par, $\text{hw} \leq N$
\Ensure Polinômio com exatamente $\text{hw}$ coeficientes não nulos:
  $\text{hw}/2$ em $+1$ e $\text{hw}/2$ em $q-1$
\State $\text{idx}[i] \gets i$ para $i = 0, \ldots, N-1$
\For{$i \gets 0$ \textbf{até} $\text{hw}-1$}
  \State $j \gets \mathcal{U}\{i, \ldots, N-1\}$
    \Comment{Inteiro uniforme: Fisher--Yates parcial}
  \State $\text{idx}[i] \leftrightarrow \text{idx}[j]$
\EndFor
\For{$i \gets 0$ \textbf{até} $\text{hw}/2 - 1$}
  \State $\text{poly}[\text{idx}[i]] \gets 1$
\EndFor
\For{$i \gets \text{hw}/2$ \textbf{até} $\text{hw}-1$}
  \State $\text{poly}[\text{idx}[i]] \gets q - 1$
\EndFor
\State \Return $\text{poly}$
\end{algorithmic}
\end{algorithm}

% =============================================================================
\chapter{Geração de Chaves}
\label{cap:keygen}
% =============================================================================

O esquema de chaves é baseado no problema RLWE~\cite{LPR2010}. A classe
\texttt{KeyGenerator} encapsula um \texttt{Sampler} e um \texttt{NttCtx} e é não
copiável. Dois objetos são gerados:

\begin{itemize}
  \item \textbf{Chave secreta} (\texttt{SecretKey}): polinômio ternário
    $s \in \{-1, 0, 1\}^N$, com $-1$ representado como $q-1$;
  \item \textbf{Chave pública} (\texttt{PublicKey}): par $(p_0, p_1) \in \Rq^2$
    satisfazendo a invariante
    \begin{equation}
      p_0 + p_1 \cdot s = e \pmod{q},
      \label{eq:pk_invariant}
    \end{equation}
    onde $e$ é um ruído gaussiano pequeno.
\end{itemize}

\begin{algorithm}[ht]
\caption{\texttt{gen\_public\_key}: geração da chave pública RLWE}
\label{alg:keygen}
\begin{algorithmic}[1]
\Require Chave secreta $s$;\; parâmetros $(N, q, \sigma)$
\Ensure Par $(p_0, p_1)$ com $p_0 + p_1 \cdot s = e$
\State $a \gets \texttt{uniform}(N, q)$
  \Comment{$p_1 = a$: polinômio uniforme em $\Rq$}
\State $e \gets \texttt{gaussian}(N, q, \sigma)$
\State $a_s \gets \texttt{poly\_ntt\_mul}(a, s)$
\For{$i \gets 0$ \textbf{até} $N-1$}
  \State $p_0[i] \gets \texttt{mod\_sub}(e[i],\; a_s[i],\; q)$
    \Comment{$p_0 = e - a \cdot s$}
\EndFor
\State \Return $(p_0,\; a)$
\end{algorithmic}
\end{algorithm}

A função \texttt{gen\_key\_bundle} expõe adicionalmente o ruído $e$ para verificação
da identidade~\eqref{eq:pk_invariant} em testes, mas este valor nunca deve ser
divulgado em produção, pois permite recuperar $s$ por resolução de sistema linear.

% =============================================================================
\chapter{Cifragem e Decifragem}
\label{cap:enc_dec}
% =============================================================================

\section{Cifragem de texto claro}
\label{sec:encrypt}

A classe \texttt{Encryptor} armazena uma cópia da chave pública, um \texttt{Sampler} e
um \texttt{NttCtx}. O método \texttt{encrypt(pt)} sorteia amostras frescas a cada
chamada, garantindo segurança semântica sob RLWE.

\begin{algorithm}[ht]
\caption{\texttt{encrypt}: cifragem de texto claro no esquema CKKS}
\label{alg:encrypt}
\begin{algorithmic}[1]
\Require Texto claro $m \in \Rq$;\; chave pública $(p_0, p_1)$;\; $(N, q, \sigma)$
\Ensure Texto cifrado $(c_0, c_1)$ com $c_0 + c_1 \cdot s \approx m$
\State $u  \gets \texttt{ternary\_uniform}(N, q)$
\State $e_0 \gets \texttt{gaussian}(N, q, \sigma)$
\State $e_1 \gets \texttt{gaussian}(N, q, \sigma)$
\State $p_1 u \gets \texttt{poly\_ntt\_mul}(p_1, u)$
\State $p_0 u \gets \texttt{poly\_ntt\_mul}(p_0, u)$
\For{$i \gets 0$ \textbf{até} $N-1$}
  \State $c_1[i] \gets \texttt{mod\_add}(p_1 u[i],\; e_1[i],\; q)$
    \Comment{$c_1 = p_1 u + e_1$}
\EndFor
\For{$i \gets 0$ \textbf{até} $N-1$}
  \State $c_0[i] \gets \texttt{mod\_add}(\texttt{mod\_add}(p_0 u[i],\; e_0[i],\; q),\; m[i],\; q)$
    \Comment{$c_0 = p_0 u + e_0 + m$}
\EndFor
\State \Return $(c_0, c_1)$
\end{algorithmic}
\end{algorithm}

\section{Análise do ruído de cifragem}
\label{sec:encrypt_noise}

A quantidade decifrada é computada explicitando a invariante~\eqref{eq:pk_invariant}:
\begin{align}
  c_0 + c_1 \cdot s
    &= (p_0 u + e_0 + m) + (p_1 u + e_1) \cdot s \notag \\
    &= u \cdot (p_0 + p_1 s) + e_0 + e_1 s + m \notag \\
    &= u \cdot e_{\text{chave}} + e_0 + e_1 s + m.
    \label{eq:enc_noise}
\end{align}
O termo de erro $\varepsilon = u \cdot e_{\text{chave}} + e_0 + e_1 s$ é um polinômio
com coeficientes delimitados em norma infinita por aproximadamente
$2N \cdot \text{tail} \cdot \sigma \ll q/2$, garantindo que a decifragem recupere $m$
com erro controlado. Após a decodificação por $\Delta$, o erro por \emph{slot} é
$|\varepsilon_{\text{slot}}| \lesssim 2N \cdot \text{tail} \cdot \sigma / \Delta$,
expressão que coincide com o ruído criptográfico individual $e$ da análise de agregação
do Capítulo~4 da dissertação.

\section{Decifragem}
\label{sec:decrypt}

A classe \texttt{Decryptor} armazena a chave secreta e um \texttt{NttCtx}. A decifragem
é uma única multiplicação polinomial seguida de adição coeficiente a coeficiente.

\begin{algorithm}[ht]
\caption{\texttt{decrypt}: decifragem de texto cifrado CKKS}
\label{alg:decrypt}
\begin{algorithmic}[1]
\Require Texto cifrado $(c_0, c_1)$;\; chave secreta $s$
\Ensure $m' = c_0 + c_1 \cdot s \approx m$ (texto claro ruidoso)
\State $c_1 s \gets \texttt{poly\_ntt\_mul}(c_1, s)$
\For{$i \gets 0$ \textbf{até} $N-1$}
  \State $m'[i] \gets \texttt{mod\_add}(c_0[i],\; c_1 s[i],\; q)$
\EndFor
\State \Return $m'$
\end{algorithmic}
\end{algorithm}

\section{Pipeline de cifragem e decifragem}
\label{sec:pipeline}

A \cref{fig:enc_pipeline} resume o fluxo completo, incluindo a etapa de avaliação
homomórfica descrita no \cref{cap:evaluator}.

\begin{figure}[ht]
\centering
\begin{tikzpicture}[
  node distance = 0.4cm and 0.9cm,
  box/.style = {draw, rounded corners=3pt, minimum width=2.3cm, minimum height=0.7cm,
                align=center, font=\small},
  io/.style  = {draw, ellipse, minimum width=1.8cm, minimum height=0.65cm,
                align=center, font=\small},
  arr/.style = {-{Stealth[length=5pt]}, thick},
  sub/.style = {font=\scriptsize, text=gray}
]
  \node[io]  (m)    {$m \in \Rq$};
  \node[box, right=0.8cm of m]   (enc)  {\texttt{encrypt}\\$(p_0,p_1,u,e_0,e_1)$};
  \node[io,  right=0.8cm of enc] (ct)   {$(c_0^{(k)}, c_1^{(k)})$};
  \node[box, right=0.8cm of ct]  (add)  {\texttt{add\_inplace}\\$k = 1,\ldots,K$};
  \node[io,  right=0.8cm of add] (cta)  {$(c_0', c_1')$};
  \node[box, right=0.8cm of cta] (dec)  {\texttt{decrypt}\\$c_0' + c_1' s$};
  \node[io,  right=0.8cm of dec] (m2)   {$m' \approx \sum m^{(k)}$};

  \draw[arr] (m)   -- (enc);
  \draw[arr] (enc) -- (ct);
  \draw[arr] (ct)  -- (add);
  \draw[arr] (add) -- (cta);
  \draw[arr] (cta) -- (dec);
  \draw[arr] (dec) -- (m2);

  \node[sub, below=0.1cm of enc] {2$\times$ NTT-mul};
  \node[sub, below=0.1cm of add] {$2N$ mod\_add};
  \node[sub, below=0.1cm of dec] {1$\times$ NTT-mul};
\end{tikzpicture}
\caption{Pipeline de cifragem, avaliação homomórfica por $K$ adições e decifragem.
As operações dominantes são as multiplicações polinomiais via NTT ($O(N\log N)$ cada).}
\label{fig:enc_pipeline}
\end{figure}

% =============================================================================
\chapter{Avaliador Homomórfico Aditivo}
\label{cap:evaluator}
% =============================================================================

A classe \texttt{Evaluator} é sem estado além dos parâmetros $(N, q)$ e todos os seus
métodos são seguros para uso concorrente sobre textos cifrados distintos. O método
\texttt{add\_inplace} realiza a adição coeficiente a coeficiente das duas componentes,
conforme o \cref{alg:add_inplace}.

\begin{algorithm}[ht]
\caption{\texttt{add\_inplace}: adição homomórfica in-place}
\label{alg:add_inplace}
\begin{algorithmic}[1]
\Require Textos cifrados $(c_0^{(1)}, c_1^{(1)})$ e $(c_0^{(2)}, c_1^{(2)})$ com mesmo $(N, q)$
\Ensure $(c_0^{(1)}, c_1^{(1)}) \gets (c_0^{(1)} + c_0^{(2)},\; c_1^{(1)} + c_1^{(2)}) \bmod q$
\For{$i \gets 0$ \textbf{até} $N-1$}
  \Comment{Laço vetorizado: \texttt{\#pragma omp simd}}
  \State $c_0^{(1)}[i] \gets \texttt{mod\_add}(c_0^{(1)}[i],\; c_0^{(2)}[i],\; q)$
\EndFor
\For{$i \gets 0$ \textbf{até} $N-1$}
  \Comment{Laço vetorizado: \texttt{\#pragma omp simd}}
  \State $c_1^{(1)}[i] \gets \texttt{mod\_add}(c_1^{(1)}[i],\; c_1^{(2)}[i],\; q)$
\EndFor
\end{algorithmic}
\end{algorithm}

\paragraph{Correção.} A propriedade homomórfica verifica-se diretamente:
\begin{align}
  (c_0^{(1)}+c_0^{(2)}) + (c_1^{(1)}+c_1^{(2)}) \cdot s
    &= (c_0^{(1)} + c_1^{(1)} s) + (c_0^{(2)} + c_1^{(2)} s) \notag \\
    &\approx m_1 + m_2.
    \label{eq:add_homo}
\end{align}

\paragraph{Crescimento do ruído.} Como demonstrado na análise de agregação do Capítulo~4
da dissertação, após $K$ adições de textos cifrados independentes, o ruído acumulado
cresce em norma RMS como $\sqrt{K} \cdot \sigma_e$, onde $\sigma_e$ é o desvio padrão
do ruído individual de cifragem. O erro por \emph{slot} após decodificação é
\begin{equation}
  \varepsilon_K \approx \frac{\sqrt{K} \cdot \sigma_e}{\Delta},
  \label{eq:noise_growth}
\end{equation}
que é a mesma expressão $\sqrt{K} \cdot e / \Delta$ da equação de erro agregado da
dissertação.

% =============================================================================
\chapter{Codificação Canônica CKKS}
\label{cap:encoder}
% =============================================================================

O codificador converte vetores de números reais ou complexos em polinômios de $\Rq$ e
vice-versa, por meio do mapa de imersão canônica $\sigma$~\cite{Cheon2017}.

\section{Mapa de imersão canônica}
\label{sec:canonical_embedding}

Dado um polinômio $f \in \mathbb{Z}[X]/(X^N+1)$, a imersão canônica avalia $f$ nas $N$
raízes primitivas $2N$-ésimas da unidade:
\begin{equation}
  \omega_j = e^{i\pi(2j+1)/N}, \quad j = 0, 1, \ldots, N-1.
  \label{eq:roots}
\end{equation}
Como $\omega_{N-1-j} = \overline{\omega_j}$, a imagem $\sigma(f)$ exibe simetria
hermitiana, de modo que apenas $N/2$ valores são independentes, definindo os
$N/2$ \emph{slots} do esquema CKKS.

\section{Codificação ($\sigma^{-1}$)}
\label{sec:encode}

O \cref{alg:encode} recebe um vetor $z \in \mathbb{C}^{N/2}$, aplica a imersão inversa
e quantiza os coeficientes com a escala $\Delta$.

\begin{algorithm}[ht]
\caption{\texttt{encode}: codificação canônica CKKS}
\label{alg:encode}
\begin{algorithmic}[1]
\Require $z \in \mathbb{C}^{N/2}$;\; escala $\Delta$;\; parâmetros $(N, q)$
\Ensure Polinômio $m \in \Rq$ tal que $\sigma(m/\Delta) \approx z$
\State $v[i] \gets z[i]$ para $i = 0, \ldots, N/2-1$
\State $v[N-1-i] \gets \overline{z[i]}$ para $i = 0, \ldots, N/2-1$
  \Comment{Simetria hermitiana}
\State $v \gets \texttt{FFT}(v)$
\For{$k \gets 0$ \textbf{até} $N-1$}
  \State $a_k \gets \operatorname{Re}(v[k] \cdot e^{-i\pi k/N}) / N$
    \Comment{Torção e normalização}
  \State $c \gets \lfloor \Delta \cdot a_k \rceil$
    \Comment{Escalonamento e arredondamento}
  \State $r \gets c \bmod q$;\quad \textbf{se} $r < 0$ \textbf{então} $r \gets r + q$
  \State $m[k] \gets r$
\EndFor
\State \Return $m$
\end{algorithmic}
\end{algorithm}

O erro de quantização por coeficiente é delimitado por $1/(2\Delta)$, produzindo erro
por \emph{slot} de magnitude $\sim N/(2\Delta)$, conforme a equação de erro de
codificação $|\varepsilon_{\text{enc}}| \leq 1/(2\Delta)$ do Capítulo~4 da dissertação.

\section{Decodificação ($\sigma$)}
\label{sec:decode}

A decodificação inverte o processo: centraliza os coeficientes em $(-q/2, q/2]$, divide
por $\Delta$, aplica a torção e a IFFT.

\begin{algorithm}[ht]
\caption{\texttt{decode}: decodificação canônica CKKS}
\label{alg:decode}
\begin{algorithmic}[1]
\Require Polinômio $m \in \Rq$;\; escala $\Delta$;\; parâmetros $(N, q)$
\Ensure $z \approx \sigma(m/\Delta) \in \mathbb{C}^{N/2}$
\For{$k \gets 0$ \textbf{até} $N-1$}
  \State $c \gets m[k]$;\quad \textbf{se} $c > q/2$ \textbf{então} $c \gets c - q$
    \Comment{Centraliza em $(-q/2,\, q/2]$}
  \State $b[k] \gets (c / \Delta) \cdot e^{i\pi k/N}$
    \Comment{Desnormalizção e torção}
\EndFor
\State $v \gets N \cdot \texttt{IFFT}(b)$
  \Comment{IFFT com normalização $\times N$ para desfazer o $1/N$ da FFT}
\State \Return $v[0], \ldots, v[N/2-1]$
\end{algorithmic}
\end{algorithm}

A condição $q > 2\Delta \cdot \max_j |z_j|$ deve ser satisfeita para que a
centralização produza o coeficiente correto.

% =============================================================================
\chapter{Multiplicação por Texto Claro (Referência)}
\label{cap:mul_plain}
% =============================================================================

A função \texttt{mul\_plain(ct, pt)} implementa a multiplicação de um texto cifrado
por um texto claro por meio do algoritmo escolar $O(N^2)$ no domínio de Montgomery.
Ela serve como implementação de referência para verificação de correção e não é
utilizada na produção da biblioteca, onde a multiplicação via NTT será adicionada
em versão futura.

A redução negacíclica $X^N \equiv -1$ é aplicada inline: produtos de coeficientes com
índice $i+j \geq N$ contribuem com sinal negativo para o coeficiente de índice $i+j-N$.

\begin{algorithm}[ht]
\caption{\texttt{mul\_plain}: multiplicação escolar negacíclica em domínio de Montgomery}
\label{alg:mul_plain}
\begin{algorithmic}[1]
\Require Texto cifrado $(c_0, c_1)$;\; texto claro $m$;\; contexto \texttt{MontgomeryCtx}
\Ensure Texto cifrado $(r_0, r_1)$ com $r_b = c_b \cdot m \bmod (X^N+1, q)$ para $b \in \{0,1\}$
\State Converta todos os coeficientes para forma de Montgomery: $\hat{c}_0, \hat{c}_1, \hat{m}$
\For{$b \in \{0, 1\}$}
  \For{$i \gets 0$ \textbf{até} $N-1$}
    \For{$j \gets 0$ \textbf{até} $N-1$}
      \State $p \gets \texttt{mont\_mul}(\hat{c}_b[i],\; \hat{m}[j])$
      \If{$i + j < N$}
        \State $r_b[i+j] \gets \texttt{mod\_add}(r_b[i+j],\; p,\; q)$
      \Else
        \State $r_b[i+j-N] \gets \texttt{mod\_sub}(r_b[i+j-N],\; p,\; q)$
          \Comment{$X^N \equiv -1$: sinal negativo}
      \EndIf
    \EndFor
  \EndFor
  \State Converta $r_b$ de volta: $\texttt{from\_mont}$ coeficiente a coeficiente
\EndFor
\State \Return $(r_0, r_1)$
\end{algorithmic}
\end{algorithm}

% =============================================================================
\chapter{Serialização}
\label{cap:serialization}
% =============================================================================

O módulo de serialização implementa um formato binário \emph{little-endian}
auto-descritivo para transporte em redes FL. A unidade atômica é a representação de
um polinômio \texttt{Poly}.

\section{Formato de fio}
\label{sec:wire_format}

O leiaute de bytes de um polinômio é:
\begin{equation}
  \underbrace{N\;(\texttt{u32,\;4\,B})}_{\text{grau}}
  \;\|\;
  \underbrace{q\;(\texttt{u64,\;8\,B})}_{\text{módulo}}
  \;\|\;
  \underbrace{c_0,\;c_1,\;\ldots,\;c_{N-1}\;(N \times 8\,\text{B})}_{\text{coeficientes em }
    \emph{little-endian}},
  \label{eq:poly_wire}
\end{equation}
totalizando $12 + 8N$ bytes por polinômio. Os objetos compostos são concatenações
diretas desse bloco:

\begin{itemize}
  \item \textbf{Ciphertext}: $[c_0\;\text{Poly}]\;[c_1\;\text{Poly}]$
    $= 2(12+8N)$ bytes;
  \item \textbf{PublicKey}: $[p_0\;\text{Poly}]\;[p_1\;\text{Poly}]$
    $= 2(12+8N)$ bytes;
  \item \textbf{SecretKey}: $[s\;\text{Poly}]$ $= 12+8N$ bytes.
\end{itemize}

A \cref{fig:serial_layout} ilustra o leiaute de um \texttt{Ciphertext}.

\begin{figure}[ht]
\centering
\begin{tikzpicture}[
  node distance = 0cm,
  hdr/.style = {draw, minimum width=1.5cm, minimum height=0.7cm,
                align=center, font=\scriptsize, fill=blue!15},
  dat/.style = {draw, minimum width=4.5cm, minimum height=0.7cm,
                align=center, font=\scriptsize, fill=gray!12},
  lbl/.style = {font=\footnotesize}
]
  % Poly c0
  \node[hdr] (N0)  at (0,1)    {$N$\\u32, 4B};
  \node[hdr, right=of N0] (q0) {$q$\\u64, 8B};
  \node[dat, right=of q0] (d0) {$c_0[0],\;c_0[1],\;\ldots,\;c_0[N{-}1]$\\$N\times$ u64, $8N$ B};

  % Poly c1
  \node[hdr, right=0.3cm of d0] (N1) {$N$\\u32, 4B};
  \node[hdr, right=of N1] (q1) {$q$\\u64, 8B};
  \node[dat, right=of q1] (d1) {$c_1[0],\;c_1[1],\;\ldots,\;c_1[N{-}1]$\\$N\times$ u64, $8N$ B};

  % Brace c0
  \draw[decorate, decoration={brace, amplitude=5pt, mirror}]
    (N0.south west) -- (d0.south east)
    node[midway, below=7pt, lbl] {componente $c_0$: $12+8N$ bytes};

  % Brace c1
  \draw[decorate, decoration={brace, amplitude=5pt, mirror}]
    (N1.south west) -- (d1.south east)
    node[midway, below=7pt, lbl] {componente $c_1$: $12+8N$ bytes};

  % Brace total
  \draw[decorate, decoration={brace, amplitude=8pt}]
    (N0.north west) -- (d1.north east)
    node[midway, above=9pt, lbl]
    {\texttt{Ciphertext}: $2(12+8N)$ bytes};
\end{tikzpicture}
\caption{Leiaute de bytes do \texttt{Ciphertext} no formato \emph{little-endian}
auto-descritivo. Cada polinômio carrega seu próprio $N$ e $q$, permitindo validação
completa de integridade na desserialização.}
\label{fig:serial_layout}
\end{figure}

\section{Validação na desserialização}
\label{sec:deser_validation}

A função \texttt{from\_bytes} verifica quatro condições antes de aceitar os dados:
(i) comprimento suficiente para o cabeçalho $(N, q)$; (ii) comprimento suficiente para
os $N$ coeficientes; (iii) concordância dos valores de $N$ e $q$ lidos com os esperados,
prevenindo substituição silenciosa de parâmetros; e (iv) ausência de bytes residuais após
o final esperado.

% =============================================================================
\chapter{Compilação e Vetorização}
\label{cap:build}
% =============================================================================

O núcleo \texttt{ckks\_core} é compilado como biblioteca estática com código independente
de posição (\texttt{-fPIC}), permitindo ligação tanto a executáveis quanto ao módulo de
extensão Python. A \cref{tab:flags} lista as diretivas de compilação e seu propósito.

\begin{table}[ht]
\centering
\caption{Diretivas de compilação do \texttt{ckks\_core} via CMake.}
\label{tab:flags}
\begin{tabular}{lp{9cm}}
\toprule
\textbf{Flag} & \textbf{Propósito} \\
\midrule
\texttt{-O3}             & Otimizações agressivas de código (desenrolamento de laço, \emph{inlining}) \\
\texttt{-march=native}   & Emite instruções para a microarquitetura do \emph{host}, habilitando AVX2 \\
\texttt{-ffast-math}     & Otimizações de ponto flutuante; utilizado principalmente pelo codificador CKKS \\
\texttt{-fopenmp-simd}   & Ativa \texttt{\#pragma omp simd} sem carregar o \emph{runtime} de \emph{threads} OpenMP \\
\bottomrule
\end{tabular}
\end{table}

A diretiva \texttt{\#pragma omp simd} está presente nos laços de adição de
\texttt{Poly::operator+} e \texttt{Evaluator::add\_inplace}. Nos laços internos da NTT,
a vetorização é induzida pelos qualificadores \texttt{\_\_restrict\_\_} nos ponteiros de
entrada e saída da borboleta, que informam ao compilador a ausência de aliasing.

Vale registrar uma divergência entre os dois caminhos de construção: o CMake
inclui \texttt{-ffast-math} e \texttt{-fopenmp-simd}; a construção alternativa via
\texttt{setup.py} (modo \emph{pip}) utiliza apenas \texttt{-O3 -march=native -Wall}.
Resultados numéricos do codificador podem diferir levemente entre os dois modos quando
operando próximo dos limites de precisão de ponto flutuante, pois \texttt{-ffast-math}
permite reordenações que alteram o arredondamento.

% =============================================================================
\chapter{Conclusão e Trabalhos Futuros}
\label{cap:conclusion}
% =============================================================================

Este documento apresentou a implementação de baixo nível do \emph{backend} C++ da
biblioteca CKKS customizada, organizada em torno de três pilares técnicos:
(i)~redução de Montgomery para eliminação de divisões na aritmética modular, com
custo de construção $O(1)$ amortizado sobre toda a execução;
(ii)~NTT negacíclica para multiplicação polinomial em $O(N \log N)$, com tabelas
de twiddle pré-computadas em forma de Montgomery e em ordem bit-reversa para acesso
linear à linha de \emph{cache}; e
(iii)~aritmética modular sem desvios condicionais para vetorização automática via AVX2.

As escolhas de projeto, em particular o uso de um único primo NTT em lugar de uma cadeia
RNS completa e a restrição ao regime somente-aditivo, mantiveram a complexidade de
implementação baixa sem comprometer a correção para o caso de uso alvo: agregação
federada de gradientes por adição homomórfica.

Os trabalhos futuros imediatos incluem:

\begin{enumerate}
  \item \textbf{NTT multithread:} paralelização dos laços de borboleta com
    \texttt{-fopenmp} para aproveitar múltiplos núcleos em grades $N \geq 8192$;
  \item \textbf{Multiplicação escalar via NTT:} substituição de \texttt{mul\_plain}
    por \texttt{poly\_ntt\_mul}, reduzindo a complexidade de $O(N^2)$ para $O(N \log N)$;
  \item \textbf{Cadeia RNS:} extensão para múltiplos primos $q_1, \ldots, q_L$ com
    implementação de \emph{rescaling}, habilitando profundidade multiplicativa para
    circuitos de FL mais complexos.
\end{enumerate}

% =============================================================================
\begin{thebibliography}{99}

\bibitem{Cheon2017}
J.~H. Cheon, A.~Kim, M.~Kim, e Y.~Song,
``Homomorphic encryption for arithmetic of approximate numbers,''
em \emph{Advances in Cryptology -- ASIACRYPT 2017},
Lecture Notes in Computer Science, vol.~10624,
Springer, 2017, pp.~409--437.

\bibitem{Montgomery1985}
P.~L. Montgomery,
``Modular multiplication without trial division,''
\emph{Mathematics of Computation},
vol.~44, n.~170, pp.~519--521, abr. 1985.

\bibitem{CooleyTukey1965}
J.~W. Cooley e J.~W. Tukey,
``An algorithm for the machine calculation of complex Fourier series,''
\emph{Mathematics of Computation},
vol.~19, n.~90, pp.~297--301, 1965.

\bibitem{LongaNaehrig2016}
P.~Longa e M.~Naehrig,
``Speeding up the number theoretic transform for faster ideal lattice-based cryptography,''
em \emph{Cryptology and Network Security -- CANS 2016},
Lecture Notes in Computer Science, vol.~10052,
Springer, 2016, pp.~124--139.

\bibitem{LPR2010}
V.~Lyubashevsky, C.~Peikert, e O.~Regev,
``On ideal lattices and learning with errors over rings,''
em \emph{Advances in Cryptology -- EUROCRYPT 2010},
Lecture Notes in Computer Science, vol.~6110,
Springer, 2010, pp.~1--23.

\bibitem{sealmanual}
Microsoft Research,
\emph{SEAL: Simple Encrypted Arithmetic Library -- Manual},
versão 4.1, 2023.
Disponível em: \url{https://github.com/microsoft/SEAL}

\bibitem{Pomerance1980}
C.~Pomerance, J.~L. Selfridge, e S.~S. Wagstaff~Jr.,
``The pseudoprimes to $25 \cdot 10^9$,''
\emph{Mathematics of Computation},
vol.~35, n.~151, pp.~1003--1026, 1980.

\bibitem{Matsumoto1998}
M.~Matsumoto e T.~Nishimura,
``Mersenne twister: A 623-dimensionally equidistributed uniform pseudo-random number
generator,''
\emph{ACM Transactions on Modeling and Computer Simulation},
vol.~8, n.~1, pp.~3--30, jan. 1998.

\bibitem{Knuth1997}
D.~E. Knuth,
\emph{The Art of Computer Programming, Vol.~2: Seminumerical Algorithms},
3.ª~ed., Addison-Wesley, 1997.

\end{thebibliography}

\end{document}
